\DTMsavedate{mydate}{2022-06-30}
\maketitle{Digital Signal Processing}{Electronic Engineering}{\DTMusedate{mydate}}

% ----------------------------------------------------- %
% COURSE INFO
\subsection*{Course Information}\label{course:tee3101}
\vspace{0ex}%
\noindent\parbox{0.5\textwidth}{%
    \noindent\begin{tabular}{@{}ll}
                 \textsf{Course:}          & 	Digital Signal Processing  \\
                 \textsf{Course Code:}         & 	TEE 3101    \\
                 \textsf{Credit Hours:}    & 	10
    \end{tabular}}
\vspace{2ex}\hrule\vspace{2ex}

% ----------------------------------------------------- %
% INSTRUCTOR INFO
\subsection*{Lecturer Information}
\noindent\parbox{0.5\textwidth}{%
    \noindent\begin{tabular}{@{}ll}
                 \textsf{Name:}         &    Zedekia Madumbu Nyathi         \\
% \textsf{Phone:}        &    \phone          \\
\textsf{Email:}        &    zedekia.madumbu.nyathi@nust.ac.zw \\
\textsf{Qualifications:}        & MSc in Electronic Engineering, FETC
    \end{tabular}}
\vspace{2ex}\hrule\vspace{2ex}


% ----------------------------------------------------- %
% COURSE DESCRIPTION
\subsection*{Course Synopsis}
Analysis of continuous and discrete signals and systems.
Fourier series and transforms.
Laplace transforms, Z transforms, transfer functions, analysis of stability, probabilistic convolution, impulse response and transfer functions.
\vspace{2ex} \hrule\vspace{2ex}

% ----------------------------------------------------- %
% COURSE TOPICS
\subsection*{Topics}
\begin{outline}
    \1 Analysis of continuous and discrete signals and systems
    \1 Time domain representation of continuous and discrete signals
    \1 Time domain characterisation of linear time invariant continuous and discrete time signals and systems
    \1 Digital Processing of continuous time signals
    \1 Fourier series and transform frequency domain representation of continuous and discrete time signals
    \1 Laplace transforms
    \1 Z transforms
    \1 Transfer functions, derivation of the transfer functions, Analysis of stability
\end{outline}
\vspace{2ex} \hrule\vspace{2ex}

% ----------------------------------------------------- %
% Students Assenssment
\subsection*{Assessment of Students}
\begin{outline}
    \1 Students are assessed through seven tutorials, assignments, two written in-class tests, all of which shall form the course work assessment mark.
    \1 There shall be final examination at the end of each semester.
\end{outline}
\vspace{2ex}\hrule\vspace{2ex}

% ----------------------------------------------------- %
% Grade Evaluation
\subsection*{Course Evaluation and Grading Policies}
Grades are based on:
Continuous assessment (\qty{25}{\percent}),
Final Exam (\qty{75}{\percent})
\vspace{2ex}\hrule\vspace{2ex}

% ----------------------------------------------------- %
% EQUIPMENT
% https://sites.udel.edu/computing-purchases/personal-specs/
% \subsection*{Essential Equipment}

% \vspace{2ex}\hrule\vspace{2ex}

% Policies and Other information
\subsection*{Policies and Other Information}
% Conduct
\subsection*{Key Text}
\begin{itemize}
    \item Introductory digital signal processing with computer applications by Paul A. Lynn, Wolfgang Faurst
    \item Digital signal processing by John G.Proakis,Dimitris G.Manolakis
    \item Digital signal processing: A Practical Approach by Emmanuel C. Ifeachor, Barrie W.Jervis
    \item Digital signal processing by Oppenheim A.V and Schafer R.W
\end{itemize}
\vspace{2ex}\hrule
\vspace{2ex}\hrule\vspace{2ex}